% ═══════════════════════════════════════════════════════════════════════
%  FICHE PEDAGOGIQUE DE MATHEMATIQUES — TERMINISH
%  Classe : 4ème A — Année scolaire 2026-2027 — 6 séances de 55 min
%  Thème 2 : Programme de calcul – calcul littéral
%  LEÇON 2 : CALCUL SUR LES EXPRESSIONS ALGEBRIQUES
%  Canevas : fiche "Propriété de Thalès" (3ème) — même mise en page.
%
%  Compatible Overleaf (pdflatex/xelatex) ET le moteur local :
%  packages utilisés : article, geometry, xcolor, amsmath, longtable.
%  Aucun package exotique (pas tikz/colortbl/multirow/fancyhdr).
% ═══════════════════════════════════════════════════════════════════════
\documentclass[10pt]{article}

\usepackage[a4paper,landscape,top=1.5cm,bottom=1.3cm,left=1.1cm,right=1.1cm]{geometry}
\usepackage{amsmath}
\usepackage[RGB]{xcolor}
\usepackage{longtable}

% ── Couleurs du canevas ────────────────────────────────────────────────
\definecolor{rougetitre}{RGB}{192,0,0}      % titres rouges soulignés
\definecolor{rougemaison}{RGB}{160,30,30}   % exercice de maison
\definecolor{vertentete}{RGB}{215,235,210}  % fond vert pâle (en-têtes)
\definecolor{bleuseance}{RGB}{46,110,180}   % cadres bleus des séances
\definecolor{bleufond}{RGB}{222,235,247}    % fond bleu pâle (n° séance)
\definecolor{bleutrace}{RGB}{30,80,150}     % mots "trace écrite"
\definecolor{jaunesurl}{RGB}{255,230,90}    % surlignage mots-clés
\definecolor{vertbilan}{RGB}{230,242,225}   % fond vert très pâle

\setlength{\parindent}{0pt}
\setlength{\tabcolsep}{3pt}
\renewcommand{\arraystretch}{1.25}

% ── Raccourcis de style ────────────────────────────────────────────────
\newcommand{\titreR}[1]{\underline{\textcolor{rougetitre}{\bfseries #1}}}
\newcommand{\app}{\titreR{Exercice d'application}}
\newcommand{\maison}{\underline{\textcolor{rougemaison}{\bfseries Exercice de maison}}}
\newcommand{\seance}[1]{\fcolorbox{bleuseance}{bleufond}{\textcolor{bleuseance}{\bfseries #1}}}
\newcommand{\moment}[1]{\textbf{\itshape #1}}
\newcommand{\mclef}[1]{\colorbox{jaunesurl}{#1}}
\newcommand{\prop}[1]{\par\medskip\noindent\fbox{\parbox{0.96\linewidth}{#1}}\par\medskip}
\newcommand{\rouge}[1]{\textcolor{rougetitre}{\bfseries #1}}
\newcommand{\EnT}{\colorbox{vertentete}{\makebox[\linewidth][c]{\bfseries}}} % cellule en-tête verte
\newcommand{\ieme}{\textsuperscript{ième}}
\newcommand{\iere}{\textsuperscript{ière}}
\newcommand{\ere}{\textsuperscript{ère}}

% ── En-tête et pied de page (comme le canevas) ────────────────────────
\makeatletter
\def\ps@fiche{%
  \def\@oddhead{\hfill\underline{\textcolor{rougetitre}{\bfseries FICHE PEDAGOGIQUE DE MATHEMATIQUES}}\hfill}%
  \def\@oddfoot{\hfill\textcolor{bleuseance}{\bfseries Maths 4\textsuperscript{eme} A}\hfill
                \fcolorbox{black}{vertentete}{\arabic{page}}\hfill}%
}
\makeatother
\pagestyle{fiche}

\begin{document}

% ═════════════════════════ CARTOUCHE D'IDENTIFICATION ═════════════════════════
\begin{tabular}{|p{8.6cm}|p{8.6cm}|p{8.6cm}|}
\hline
\textbf{NOM} \; : \; {\bfseries KODJONE} &
\textbf{DRE} \; : \; \dotfill &
\textbf{FICHE N\textsuperscript{o}} \; : \; {\bfseries 1} \\ \hline
\textbf{PRÉNOM} \; : \; {\bfseries Kodjo} &
\textbf{IESG} \; : \; \dotfill &
\textbf{DISCIPLINE} \; : \; {\bfseries MATHS} \\ \hline
\textbf{CONTACT} \; : \; \dotfill &
\textbf{ÉTABLISSEMENT} \; : \; \dotfill &
\textbf{CLASSE} \; : \; {\bfseries 4\ieme{} A} \\ \hline
\textbf{EFFECTIF} \; : \; \textbf{G} : \dotfill \quad
                           \textbf{F} : \dotfill \quad
                           \textbf{T} : \dotfill &
\textbf{ANNÉE SCOLAIRE} \; : \; {\bfseries 2026 - 2027} &
\textbf{SEMAINE} \; : \; {\bfseries 1} \\ \hline
\end{tabular}

\bigskip
\fcolorbox{black}{vertbilan}{%
\begin{minipage}{27.1cm}
\smallskip
{\bfseries COMPETENCE TERMINALE 2 :} \textit{Résoudre des problèmes de la vie courante faisant appel
à des programmes de calcul, à l'activité calcul littéral et à l'organisation de calculs
à l'aide d'expressions algébriques.}\\[3pt]
{\bfseries THEME 2 :} \hfill {\Large\bfseries PROGRAMME DE CALCUL – CALCUL LITTERAL} \hfill\mbox{}\\[3pt]
{\bfseries LEÇON 2 :} \hfill \fcolorbox{rougetitre}{white}{\textcolor{rougetitre}{\large\bfseries CALCUL SUR LES EXPRESSIONS ALGEBRIQUES}} \hfill\mbox{}\\[3pt]
{\bfseries NOMBRE DE SÉANCES :} \; {\bfseries 6} \qquad
{\bfseries DURÉE D'UNE SÉANCE :} \; {\bfseries 55 min}\\[3pt]
{\bfseries SUPPORTS DIDACTIQUES PRINCIPAUX :} \; Enoncé de la situation d'apprentissage ; CIAM 4\ieme ; CARGO 4\ieme ; Programme et son guide\\[3pt]
{\bfseries MATERIEL DIDACTIQUE ORDINAIRE :} \; Régie ; craie ; règle ; cahiers \qquad
{\bfseries MATERIEL DIDACTIQUE ADAPTÉ :} \; Fiche d'activités ; expressions découpées (étiquettes)\\[3pt]
{\bfseries CONSIGNE POUR LE PROFESSEUR DE SUIVI :} \; Vérifier le respect des moments didactiques et la trace écrite de la leçon.\\[3pt]
{\bfseries PREREQUIS :} \; Priorités opératoires ; programmes de calcul ; vocabulaire (somme, différence, produit, quotient, terme, facteur) ;
propriétés de la multiplication (commutativité, associativité, distributivité sur l'addition).
\smallskip
\end{minipage}}

\bigskip
% ═════════════════════════ TABLEAU CAPACITES / CONTENUS ═════════════════════════
\begin{tabular}{|p{9.4cm}|p{17.3cm}|}
\hline
\multicolumn{1}{|c|}{\colorbox{vertentete}{\makebox[9.2cm][c]{\bfseries Capacités}}} &
\multicolumn{1}{c|}{\colorbox{vertentete}{\makebox[17.1cm][c]{\bfseries Contenus}}} \\ \hline
Reconnaître une expression algébrique et son vocabulaire &
\small Expressions algébriques : \mclef{variable}, \mclef{constante}, \mclef{terme}, \mclef{coefficient} ; traduction d'un programme de calcul par une expression algébrique. \\ \hline
Réduire une expression algébrique &
\small Réduction : \small termes semblables ; regroupement ; $ax+bx=(a+b)\,x$ ; ordre décroissant des puissances. \\ \hline
Développer une expression algébrique &
\small Développement : $k\,(a+b)=ka+kb$ ; \; $(a+b)\,(c+d)=ac+ad+bc+bd$. \\ \hline
Factoriser une expression algébrique &
\small Factorisation par \small facteur commun : $ka+kb=k\,(a+b)$. \\ \hline
Calculer la valeur numérique d'une expression algébrique &
\small Substitution : remplacer la variable \small par la valeur donnée, en respectant les priorités opératoires. \\ \hline
Justifier l'égalité de deux écritures d'une même expression &
\small Propriétés de la distributivité \small (développer / factoriser) ; périmètres, aires, programmes de calcul. \\ \hline
\end{tabular}

\bigskip
{\bfseries Stratégie pédagogique}

\begin{itemize}
  \item[] — \; Travail individuel
  \item[] — \; Travail en groupe
  \item[] — \; Les discussions
\end{itemize}

\bigskip
% ═════════════════════════ DEROULEMENT : TABLEAU A 4 COLONNES ═════════════════════════
\begin{longtable}{|p{2.35cm}|p{5.35cm}|p{4.55cm}|p{13.85cm}|}
% ── en-tête répété sur chaque page (comme le canevas) ──
\hline
\multicolumn{1}{|c|}{\colorbox{vertentete}{\parbox{2.15cm}{\centering\small\bfseries Moment didactique}}} &
\multicolumn{1}{c|}{\colorbox{vertentete}{\makebox[5.2cm][c]{\small\bfseries Activité de l'enseignant}}} &
\multicolumn{1}{c|}{\colorbox{vertentete}{\makebox[4.4cm][c]{\small\bfseries Activité de l'élève}}} &
\multicolumn{1}{c|}{\colorbox{vertentete}{\makebox[13.7cm][c]{\small\bfseries Trace écrite}}} \\
\hline
\endfirsthead
\hline
\multicolumn{1}{|c|}{\colorbox{vertentete}{\parbox{2.15cm}{\centering\small\bfseries Moment didactique}}} &
\multicolumn{1}{c|}{\colorbox{vertentete}{\makebox[5.2cm][c]{\small\bfseries Activité de l'enseignant}}} &
\multicolumn{1}{c|}{\colorbox{vertentete}{\makebox[4.4cm][c]{\small\bfseries Activité de l'élève}}} &
\multicolumn{1}{c|}{\colorbox{vertentete}{\makebox[13.7cm][c]{\small\bfseries Trace écrite}}} \\
\hline
\endhead
% ═════════════════════ SÉANCE 1 ═════════════════════
\seance{SEANCE 1} \newline \moment{Prérequis} &
\titreR{Exercice} : Consigne d'effectuer au tableau les calculs
$A=3\times5+4\times2$, \; $B=2\times(7+3)$, puis de dire, en français, ce qu'on a fait
(« j'ai multiplié… puis j'ai ajouté… »). Corrige et fait le point. &
\textbf{Résolution} : les élèves effectuent les calculs en respectant les priorités
et décrivent oralement le programme de calcul suivi. &
\titreR{Exercice prérequis} \par $A=3\times5+4\times2=15+8=23$ \par
$B=2\times(7+3)=2\times10=20$ \par
\mclef{Programme de calcul} : choisir un nombre, le multiplier par $3$, ajouter $4$. \\ \hline

\moment{Présentation de la situation-problème (5 min)} &
Recopie la situation au tableau et la lit à haute voix. &
Recopient la situation dans le cahier et la lisent. &
\titreR{Situation d'apprentissage} \par
\textit{Pour agrandir son poulailler, M. ABALO veut construire un nouveau parc
rectangulaire de longueur $(3x+2)$~m et de largeur $(x+4)$~m. Pour connaître la longueur
de grillage nécessaire, il demande à ses trois enfants de calculer le périmètre $P$
du rectangle. Le premier écrit : $P=2\,(3x+2)+2\,(x+4)$ ; le deuxième écrit :
$P=(3x+2)+(3x+2)+(x+4)+(x+4)$ ; la troisième réduit et écrit : $P=8x+12$.} \par
\begin{center}
\setlength{\unitlength}{1mm}%
\begin{picture}(84,26)
  \put(2,2){\framebox(56,17){}}
  \put(15,21){\small longueur $=(3x+2)$ m}
  \put(21,8){\small parc rectangulaire}
  \put(62,12){\small largeur}
  \put(62,8){\small $=(x+4)$ m}
\end{picture}
\end{center}
\textbf{Problème :} les trois écritures donnent-elles le même périmètre ? Laquelle est
la plus simple ? \\ \hline

\moment{Appropriation de la situation-problème, compréhension de la tâche et de
l'organisation du travail (10 min)} &
Explique les mots difficiles, pose des questions de compréhension et répond aux
interrogations des élèves. &
Exploquent la situation, posent des questions et relèvent les mots difficiles. &
\rouge{Consignes :} \par \textbf{1)} Traduire le programme de calcul de chacun des trois enfants. \par
\textbf{2)} Montrer que les trois écritures donnent le même périmètre. \par
\textbf{3)} Donner l'écriture la plus simple de $P$, puis calculer ce périmètre pour $x=2$~m. \\ \hline

\moment{Travail individuel (5 min) et travail en groupe (15 min)} &
Vérifie les productions des élèves, fait des remarques à haute voix, circule et guide
les groupes sans donner la solution. &
\textbf{Résolution} : cherchent individuellement puis en groupe ; confrontent leurs
écritures ; rédigent une production. &
\titreR{Recherche} \par
\textbf{1)} $1^{\text{er}}$ enfant : choisir $x$, calculer $3x+2$, le doubler, puis
calculer $x+4$, le doubler, et ajouter. \par
$2^{\text{e}}$ enfant : additionner les quatre côtés. \par
$3^{\text{e}}$ enfant : regrouper les termes semblables. \par
\textbf{2)} $2\,(3x+2)+2\,(x+4)=6x+4+2x+8=(6x+2x)+(4+8)=8x+12$ ; \par
$(3x+2)+(3x+2)+(x+4)+(x+4)=(3x+3x+x+x)+(2+2+4+4)=8x+12$. \par
\textbf{3)} Écriture la plus simple : $P=8x+12$ ; pour $x=2$ : $P=8\times2+12=28$~m. \\ \hline

\moment{Synthèse et bilan du travail (20 min)} &
Apporte des amendements à la présentation des élèves puis valide la production. &
Présentent leurs travaux au tableau ; apportent des amendements à la présentation. &
\titreR{Synthèse} \par
\textbf{Données :} longueur $=3x+2$ ; largeur $=x+4$ ; $P$ le périmètre. \par
\textbf{Calculs :} $P=L+l+L+l=2\,(3x+2)+2\,(x+4)$ \par
\hfill $=6x+4+2x+8=(6x+2x)+(4+8)=8x+12$. \par
Pour $x=2$~m : $P=8\times2+12=28$~m. \par
\textbf{Conclusion :} les trois écritures sont égales ; la plus simple est $P=8x+12$.
On dit qu'on l'a \mclef{réduite}. \\ \hline
\end{longtable}

\begin{center}
\underline{\textcolor{rougetitre}{\large\bfseries Institutionalisation \; : \; (trace écrite de la leçon par le professeur)}}
\end{center}

% ═════════════════════ SÉANCE 2 ═════════════════════
\begin{longtable}{|p{2.35cm}|p{5.35cm}|p{4.55cm}|p{13.85cm}|}
\hline
\seance{SEANCE 2} \newline \moment{Correction de l'exercice de maison} &
Fait passer des élèves au tableau, corrige et précise les erreurs. &
\textbf{Résolution} : deux élèves traitent l'exercice au tableau ; les autres corrigent. &
\titreR{Correction de l'exercice de maison} \par
Réduire : $A=7x+3-2x+5=(7x-2x)+(3+5)=5x+8$ ; \par
$B=5a+2b-a+4b=(5a-a)+(2b+4b)=4a+6b$. \\ \hline

\moment{Leçon du jour (15 min)} &
Présente le vocabulaire à partir d'exemples d'expressions, fait remarquer les termes
semblables, puis énonce la règle de réduction. &
Écoutent, répondent aux questions, identifient termes, coefficients et variables ;
recopient la leçon. &
\titreR{Trace écrite} \par
\textbf{I) Présentation et vocabulaire} \par
\textbf{1) Définitions et vocabulaire} \par
Une \mclef{expression algébrique} est une écriture dans laquelle des nombres et des
lettres sont reliés par des opérations. \par
\textit{Exemples :} $A=3x+5$ ; \; $B=2x^2-4x+7$ ; \; $C=(2x+1)\,(x-3)$. \par
— la lettre $x$ s'appelle la \mclef{variable} ; \par
— $3x$ est un \mclef{terme} de coefficient $3$ ; \par
— $5$ est une \mclef{constante}. \par
\textit{Exercice d'application :} dans $B=2x^2-4x+7$, donner la variable, les termes,
leurs coefficients et la constante. \par
\textit{Réponse :} variable $x$ ; termes $2x^2$, $-4x$, $7$ ; coefficients $2$, $-4$, $7$ ;
constante $7$. \\ \hline

\moment{Exercice d'application (15 min)} &
Propose l'exercice d'application, fait passer deux élèves au tableau et corrige. &
\textbf{Résolution} : résolvent dans les cahiers ; deux élèves passent au tableau. &
\prop{\textbf{\textcolor{rougetitre}{Propriété (réduction)}} \par
Réduire une expression algébrique, c'est regrouper ses \textbf{termes semblables}
(termes en $x^2$ ensemble, termes en $x$ ensemble, constantes ensemble) :
\[ a\,x+b\,x=(a+b)\,x . \]}
\app \par
Réduire : \; $R=8x+5-3x+7$ ; \; $S=4x^2+2x-5x^2+x$ ; \; $T=3a+2b-5+a-b$. \par
\textbf{Réponses :} $R=5x+12$ ; \; $S=-x^2+3x$ ; \; $T=4a+b-5$. \\ \hline

\moment{Exercice de maison (5 min)} &
Donne l'exercice à faire à la maison et annonce la suite. &
Notent l'exercice de maison. &
\maison \par
Réduire les expressions : \; $E=12y-8+3y+8$ ; \;
$F=7x^2+3x-4+x^2-3x$ ; \; $G=2(3x+1)+x-4$. \par
\textit{(CIAM 4\ieme{} — série d'exercices proposés par le professeur)} \\ \hline
\end{longtable}

% ═════════════════════ SÉANCE 3 ═════════════════════
\begin{longtable}{|p{2.35cm}|p{5.35cm}|p{4.55cm}|p{13.85cm}|}
\hline
\seance{SEANCE 3} \newline \moment{Correction de l'exercice de maison} &
Fait corriger l'exercice de maison au tableau. &
\textbf{Résolution} : correction au tableau. &
\titreR{Correction de l'exercice de maison} \par
$E=(12y+3y)+(-8+8)=15y$ ; \par
$F=(7x^2+x^2)+(3x-3x)-4=8x^2-4$ ; \par
$G=6x+2+x-4=7x-2$. \\ \hline

\moment{Leçon du jour (15 min)} &
Fait découvrir la distributivité à partir de l'aire d'un rectangle de côtés $k$ et
$(a+b)$, puis énonce la propriété de développement. &
Observent la figure, énoncent la propriété, recopient la leçon. &
\titreR{Trace écrite} \par
\textbf{2) Développement d'une expression algébrique} \par
\begin{center}
\setlength{\unitlength}{1mm}%
\begin{picture}(88,28)
  \put(2,2){\framebox(40,20){}}
  \put(20.6,2){\rule{0.35mm}{20mm}}
  \put(10,24){\small $a$}
  \put(31,24){\small $b$}
  \put(-1.5,10){\small $k$}
  \put(46,22){\small aire $=k\,(a+b)$}
  \put(46,15){\small $=$ aire $(a\times k)$ $+$ aire $(b\times k)$}
  \put(46,8){\small $=ka+kb$}
\end{picture}
\end{center}
\prop{\textbf{\textcolor{rougetitre}{Propriété (développement)}} \par
Pour tous nombres $k$, $a$, $b$, $c$, $d$ :
\[ k\,(a+b)=ka+kb \qquad ; \qquad (a+b)\,(c+d)=ac+ad+bc+bd . \]}
\textit{Exemples :} \par
$3\,(2x+5)=6x+15$ ; \; $-2\,(4x-3)=-8x+6$ ; \par
$(x+3)\,(x+2)=x^2+2x+3x+6=x^2+5x+6$. \\ \hline

\moment{Exercice d'application (15 min)} &
Propose l'exercice, fait passer deux élèves au tableau, corrige. &
\textbf{Résolution} : cherchent puis passent au tableau. &
\app \par
Développer puis réduire : \; $A=5\,(2x+3)$ ; \; $B=(x+4)\,(x+1)$ ; \;
$C=3\,(x-2)+2\,(2x+5)$. \par
\textbf{Réponses :} $A=10x+15$ ; \; $B=x^2+5x+4$ ; \;
$C=3x-6+4x+10=7x+4$. \\ \hline

\moment{Exercice de maison (5 min)} &
Donne l'exercice de maison. &
Notent l'exercice. &
\maison \par
Développer puis réduire : \; $D=4\,(3x-5)$ ; \; $E=(2x+1)\,(x+3)$ ; \;
$F=5\,(x+2)-3\,(x-1)$. \\ \hline
\end{longtable}

% ═════════════════════ SÉANCE 4 ═════════════════════
\begin{longtable}{|p{2.35cm}|p{5.35cm}|p{4.55cm}|p{13.85cm}|}
\hline
\seance{SEANCE 4} \newline \moment{Correction de l'exercice de maison} &
Fait corriger l'exercice de maison. &
\textbf{Résolution} : correction au tableau. &
\titreR{Correction de l'exercice de maison} \par
$D=12x-20$ ; \; $E=2x^2+6x+x+3=2x^2+7x+3$ ; \par
$F=5x+10-3x+3=2x+13$. \\ \hline

\moment{Leçon du jour (15 min)} &
Fait remarquer, sur des exemples, qu'on peut passer de $ka+kb$ à $k(a+b)$ :
c'est la factorisation. Énonce la propriété. &
Répondent, identifient le facteur commun, recopient la leçon. &
\titreR{Trace écrite} \par
\textbf{3) Factorisation d'une expression algébrique} \par
\prop{\textbf{\textcolor{rougetitre}{Propriété (factorisation)}} \par
Factoriser, c'est transformer une somme en produit en recherchant le
\textbf{facteur commun} :
\[ k\,a+k\,b=k\,(a+b) . \]}
\textit{Exemples :} \par
$6x+9=3\,(2x+3)$ ; \; $5x^2-15x=5x\,(x-3)$ ; \par
$(2x+1)\,(x-3)+(2x+1)\,(x+5)=(2x+1)\,(2x+2)$. \\ \hline

\moment{Exercice d'application (15 min)} &
Propose l'exercice, fait passer deux élèves au tableau, corrige. &
\textbf{Résolution} : cherchent puis passent au tableau. &
\app \par
Factoriser : \; $A=7x+21$ ; \; $B=12x^2-8x$ ; \; $C=4\,(x+1)+(x+1)\,(3x-2)$. \par
\textbf{Réponses :} $A=7\,(x+3)$ ; \; $B=4x\,(3x-2)$ ; \;
$C=(x+1)\,(4+3x-2)=(x+1)\,(3x+2)$. \\ \hline

\moment{Exercice de maison (5 min)} &
Donne l'exercice de maison. &
Notent l'exercice. &
\maison \par
Factoriser : \; $D=9x-12$ ; \; $E=20x^2+5x$ ; \; $F=6x\,(x-2)-3\,(x-2)$. \\ \hline
\end{longtable}

% ═════════════════════ SÉANCE 5 ═════════════════════
\begin{longtable}{|p{2.35cm}|p{5.35cm}|p{4.55cm}|p{13.85cm}|}
\hline
\seance{SEANCE 5} \newline \moment{Correction de l'exercice de maison} &
Fait corriger l'exercice de maison. &
\textbf{Résolution} : correction au tableau. &
\titreR{Correction de l'exercice de maison} \par
$D=3\,(3x-4)$ ; \; $E=5x\,(4x+1)$ ; \; $F=3\,(x-2)\,(2x-1)$. \\ \hline

\moment{Leçon du jour (15 min)} &
Présente la méthode de calcul de la valeur numérique d'une expression, insiste sur
les priorités opératoires et les substitutions avec parenthèses. &
Écoutent, exécutent la méthode sur un exemple, recopient la leçon. &
\titreR{Trace écrite} \par
\textbf{4) Valeur numérique d'une expression algébrique} \par
Calculer la \mclef{valeur numérique} d'une expression pour $x=a$, c'est remplacer
$x$ par $a$ (entre parenthèses) puis effectuer les calculs en respectant les priorités
opératoires. \par
\textit{Exemple :} pour $x=2$, \; $A=3x^2-2x+1$ devient \par
$A=3\times(2)^2-2\times(2)+1=3\times4-4+1=12-4+1=9$. \\ \hline

\moment{Exercice d'application (15 min)} &
Propose l'exercice, fait passer deux élèves au tableau, corrige. &
\textbf{Résolution} : cherchent puis passent au tableau. &
\app \par
Soit $B=2x^2-3x+4$ et $C=(x+1)\,(x-2)$. \par
\textbf{1)} Calculer $B$ pour $x=-1$ puis pour $x=3$. \par
\textbf{2)} Calculer $C$ pour $x=5$ ; vérifier en développant $C$ d'abord. \par
\textbf{Réponses :} $B(-1)=2+3+4=9$ ; \; $B(3)=18-9+4=13$ ; \par
$C(x)=x^2-x-2$ donc $C(5)=25-5-2=18$. \\ \hline

\moment{Exercice de maison (5 min)} &
Donne l'exercice de maison. &
Notent l'exercice. &
\maison \par
Soit $D=3x^2+2x-5$ et $E=(2x-1)\,(x+3)$. \par
\textbf{1)} Calculer $D$ pour $x=-2$ et pour $x=1{,}5$. \par
\textbf{2)} Développer $E$ puis calculer $E$ pour $x=2$. \\ \hline
\end{longtable}

% ═════════════════════ SÉANCE 6 ═════════════════════
\begin{longtable}{|p{2.35cm}|p{5.35cm}|p{4.55cm}|p{13.85cm}|}
\hline
\seance{SEANCE 6} \newline \moment{Correction de l'exercice de maison} &
Fait corriger l'exercice de maison. &
\textbf{Résolution} : correction au tableau. &
\titreR{Correction de l'exercice de maison} \par
$D(-2)=12-4-5=3$ ; \; $D(1{,}5)=6{,}75+3-5=4{,}75$ ; \par
$E=2x^2+6x-x-3=2x^2+5x-3$ ; \; $E(2)=8+10-3=15$. \\ \hline

\moment{Exercices de consolidation (30 min)} &
Propose des exercices mêlant réduction, développement, factorisation et valeur
numérique ; fait passer des élèves au tableau ; corrige. &
\textbf{Résolution} : cherchent individuellement, puis passent au tableau. &
\titreR{Exercices proposés (consolidation)} \par
\textbf{1)} Réduire : \; $H=5x^2-2x+3x^2+x-7$. \par
\textbf{2)} Développer puis réduire : \; $I=2\,(3x-1)+(x+4)\,(x-2)$. \par
\textbf{3)} Factoriser : \; $J=8x^2+12x$ ; \; $K=5\,(x-3)-(x-3)\,(2x+1)$. \par
\textbf{4)} Un programme de calcul : \textit{choisir un nombre, le multiplier par $4$,
ajouter $3$, multiplier le résultat par $2$.} \par
a) Traduire par une expression algébrique. \;
b) Calculer le résultat pour $5$. \;
c) Retrouve-t-on $46$ avec l'écriture réduite ? \par
\textbf{Réponses :} $H=8x^2-x-7$ ; \; $I=6x-2+x^2+2x-8=x^2+8x-10$ ; \par
$J=4x\,(2x+3)$ ; \; $K=(x-3)\,(5-2x-4)=(x-3)\,(1-2x)$ ; \par
4) a) $2\,(4x+3)$ ; \; b) $2\times(20+3)=46$ ; \; c) écriture réduite : $8x+6$,
pour $5$ : $46$. ✓ \\ \hline

\moment{Bilan de la leçon et exercice de maison (10 min)} &
Fait reformuler ce qu'il faut retenir ; donne l'exercice de synthèse à la maison ;
annonce la leçon suivante. &
Réforment la synthèse, notent l'exercice. &
\titreR{Synthèse (à retenir)} \par
— \textbf{Réduire} : regrouper les termes semblables ($ax+bx=(a+b)x$). \par
— \textbf{Développer} : $k(a+b)=ka+kb$ ; $(a+b)(c+d)=ac+ad+bc+bd$. \par
— \textbf{Factoriser} : $ka+kb=k(a+b)$ (facteur commun). \par
— \textbf{Valeur numérique} : remplacer la variable par sa valeur. \par
\maison \par
Un champ rectangulaire a pour longueur $(2x+5)$~m et pour largeur $(x+3)$~m. \par
\textbf{1)} Exprimer le périmètre $P$ et l'aire $S$ du champ en fonction de $x$
(écritures réduites). \par
\textbf{2)} Calculer $P$ et $S$ pour $x=4$~m. \par
\textbf{Leçon suivante :} identités remarquables (développement du carré d'une somme). \\ \hline
\end{longtable}

\begin{center}
\fcolorbox{black}{vertbilan}{\begin{minipage}{24cm}\smallskip
\centering TERMINISH — l'univers de Mathématiques \quad $\cdot$ \quad
Fiche pédagogique — Maths 4\ieme{} A — Calcul sur les expressions algébriques \quad $\cdot$ \quad
Année scolaire 2026--2027\smallskip
\end{minipage}}
\end{center}

\end{document}
